\section{Continuous maps and sheaves}
\label{section-presheaves-functorial}

\noindent
Let $f : X \to Y$ be a continuous map of topological spaces.
We will define the pushforward and pullback functors for
presheaves and sheaves.

\medskip\noindent
Let $\mathcal{F}$ be a presheaf of sets on $X$. We define the
{\it pushforward} of $\mathcal{F}$ by the rule
$$
f_*\mathcal{F}(V) = \mathcal{F}(f^{-1}(V))
$$
for any open $V \subset Y$.
Given $V_1 \subset V_2 \subset Y$ open the restriction map
is given by the commutativity of the diagram
$$
\xymatrix{
f_*\mathcal{F}(V_2) \ar[d] \ar@{=}[r] &
\mathcal{F}(f^{-1}(V_2)) \ar[d]^{\text{restriction for }\mathcal{F}} \\
f_*\mathcal{F}(V_1) \ar@{=}[r] &
\mathcal{F}(f^{-1}(V_1))
}
$$
It is clear that this defines a presheaf of sets. The construction
is clearly functorial in the presheaf $\mathcal{F}$ and hence
we obtain a functor
$$
f_* : \textit{PSh}(X) \longrightarrow \textit{PSh}(Y).
$$

\begin{lemma}
\label{lemma-pushforward-sheaf}
Let $f : X \to Y$ be a continuous map.
Let $\mathcal{F}$ be a sheaf of sets on $X$.
Then $f_*\mathcal{F}$ is a sheaf on $Y$.
\end{lemma}

\begin{proof}
This immediately follows from the fact that
if $V = \bigcup V_j$ is an open covering in $Y$,
then $f^{-1}(V) = \bigcup f^{-1}(V_j)$ is an open covering in $X$.
\end{proof}

\noindent
As a consequence we obtain a functor
$$
f_* : \Sh(X) \longrightarrow \Sh(Y).
$$
This is compatible with composition in the following
strong sense.

\begin{lemma}
\label{lemma-pushforward-composition}
Let $f : X \to Y$ and $g : Y \to Z$ be continuous maps
of topological spaces. The functors $(g \circ f)_*$
and $g_* \circ f_*$ are equal (on both presheaves
and sheaves of sets).
\end{lemma}

\begin{proof}
This is because $(g \circ f)_*\mathcal{F}(W) =
\mathcal{F}((g \circ f)^{-1}W)$ and
$(g_* \circ f_*)\mathcal{F}(W) = \mathcal{F}(f^{-1} g^{-1} W)$
and $(g \circ f)^{-1}W = f^{-1} g^{-1} W$.
\end{proof}

\noindent
Let $\mathcal{G}$ be a presheaf of sets on $Y$.
The {\it pullback presheaf} $f_p\mathcal{G}$
of a given presheaf $\mathcal{G}$ is defined as the left adjoint
of the pushforward $f_*$ on presheaves. In other words it
should be a presheaf $f_p \mathcal{G}$ on $X$ such that
$$
\Mor_{\textit{PSh}(X)}(f_p\mathcal{G}, \mathcal{F})
=
\Mor_{\textit{PSh}(Y)}(\mathcal{G}, f_*\mathcal{F}).
$$
By the Yoneda lemma this determines the pullback uniquely.
It turns out that it actually exists.

\begin{lemma}
\label{lemma-pullback-presheaves}
Let $f : X \to Y$ be a continuous map.
There exists a functor
$f_p : \textit{PSh}(Y) \to \textit{PSh}(X)$
which is left adjoint to $f_*$. For a presheaf
$\mathcal{G}$ it is determined by the rule
$$
f_p\mathcal{G}(U) = \colim_{f(U) \subset V} \mathcal{G}(V)
$$
where the colimit is over the collection of open neighbourhoods
$V$ of $f(U)$ in $Y$. The colimits are over
directed partially ordered sets.
(The restriction mappings of $f_p\mathcal{G}$ are explained in the proof.)
\end{lemma}

\begin{proof}
The colimit is over the partially ordered set consisting of open subsets
$V \subset Y$ which contain $f(U)$ with ordering by reverse inclusion.
This is a directed partially ordered set, since if $V, V'$ are in it then
so is $V \cap V'$. Furthermore, if
$U_1 \subset U_2$, then every open neighbourhood of $f(U_2)$
is an open neighbourhood of $f(U_1)$. Hence the system defining
$f_p\mathcal{G}(U_2)$ is a subsystem of the one defining
$f_p\mathcal{G}(U_1)$ and we obtain a restriction map (for
example by applying the generalities in Categories,
Lemma \ref{categories-lemma-functorial-colimit}).

\medskip\noindent
Note that the construction of the colimit is clearly functorial
in $\mathcal{G}$, and similarly for the restriction mappings.
Hence we have defined $f_p$ as a functor.

\medskip\noindent
A small useful remark is that there exists
a canonical map $\mathcal{G}(U) \to f_p\mathcal{G}(f^{-1}(U))$,
because the system of open neighbourhoods
of $f(f^{-1}(U))$ contains the element $U$. This is compatible
with restriction mappings. In other words, there is a
canonical map $i_\mathcal{G} : \mathcal{G} \to f_* f_p \mathcal{G}$.

\medskip\noindent
Let $\mathcal{F}$ be a presheaf of sets on $X$.
Suppose that $\psi : f_p\mathcal{G} \to \mathcal{F}$
is a map of presheaves of sets. The corresponding map
$\mathcal{G} \to f_*\mathcal{F}$ is the map
$f_*\psi \circ i_\mathcal{G} :
\mathcal{G} \to f_* f_p \mathcal{G} \to f_* \mathcal{F}$.

\medskip\noindent
Another small useful remark is that there exists a
canonical map $c_\mathcal{F} : f_p f_* \mathcal{F} \to \mathcal{F}$.
Namely, let $U \subset X$ open.
For every open neighbourhood $V \supset f(U)$ in $Y$
there exists a map
$f_*\mathcal{F}(V) = \mathcal{F}(f^{-1}(V))\to \mathcal{F}(U)$,
namely the restriction map on $\mathcal{F}$. And this is compatible
with the restriction mappings between values of $\mathcal{F}$
on $f^{-1}$ of varying opens containing $f(U)$. Thus we obtain
a canonical map $f_p f_* \mathcal{F}(U) \to \mathcal{F}(U)$.
Another trivial verification shows that these maps are compatible
with restriction maps and define a map $c_\mathcal{F}$
of presheaves of sets.

\medskip\noindent
Suppose that $\varphi : \mathcal{G} \to f_*\mathcal{F}$
is a map of presheaves of sets. Consider $f_p\varphi :
f_p \mathcal{G} \to f_p f_* \mathcal{F}$.
Postcomposing with $c_\mathcal{F}$ gives the desired map
$c_\mathcal{F} \circ f_p\varphi : f_p\mathcal{G} \to \mathcal{F}$.
We omit the verification that this construction is inverse
to the construction in the other direction given above.
\end{proof}

\begin{lemma}
\label{lemma-stalk-pullback-presheaf}
Let $f : X \to Y$ be a continuous map.
Let $x \in X$. Let $\mathcal{G}$ be a presheaf of sets on $Y$.
There is a canonical bijection of stalks
$(f_p\mathcal{G})_x = \mathcal{G}_{f(x)}$.
\end{lemma}

\begin{proof}
This you can see as follows
\begin{eqnarray*}
(f_p\mathcal{G})_x
& = &
\colim_{x \in U} f_p\mathcal{G}(U) \\
& = &
\colim_{x \in U} \colim_{f(U) \subset V} \mathcal{G}(V) \\
& = &
\colim_{f(x) \in V} \mathcal{G}(V) \\
& = &
\mathcal{G}_{f(x)}
\end{eqnarray*}
Here we have used
Categories, Lemma \ref{categories-lemma-colimits-commute},
and the fact that any $V$ open in $Y$ containing $f(x)$
occurs in the third description above. Details omitted.
\end{proof}

\noindent
Let $\mathcal{G}$ be a sheaf of sets on $Y$.
The {\it pullback sheaf} $f^{-1}\mathcal{G}$ is defined
by the formula
$$
f^{-1}\mathcal{G} = (f_p\mathcal{G})^\# .
$$
The pullback $f^{-1}$ is a left adjoint of
pushforward on sheaves. In other words,
$$
\Mor_{\Sh(X)}(f^{-1}\mathcal{G}, \mathcal{F})
=
\Mor_{\Sh(Y)}(\mathcal{G}, f_*\mathcal{F}).
$$
Namely, we have
\begin{eqnarray*}
\Mor_{\Sh(X)}(f^{-1}\mathcal{G}, \mathcal{F})
& = &
\Mor_{\textit{PSh}(X)}(f_p\mathcal{G}, \mathcal{F}) \\
& = &
\Mor_{\textit{PSh}(Y)}(\mathcal{G}, f_*\mathcal{F}) \\
& = &
\Mor_{\Sh(Y)}(\mathcal{G}, f_*\mathcal{F})
\end{eqnarray*}
For the first equality we use that sheafification is a left adjoint
to the inclusion of sheaves in presheaves. For the second equality
we use that $f_p$ is a left adjoint to $f_*$ on presheaves. We will return
to this statement in the proof of Lemma \ref{lemma-f-map}.

\begin{lemma}
\label{lemma-stalk-pullback}
Let $x \in X$. Let $\mathcal{G}$ be a sheaf of sets on $Y$.
There is a canonical bijection of stalks
$(f^{-1}\mathcal{G})_x = \mathcal{G}_{f(x)}$.
\end{lemma}

\begin{proof}
This is a combination of Lemmas \ref{lemma-stalk-sheafification}
and \ref{lemma-stalk-pullback-presheaf}.
\end{proof}

\begin{lemma}
\label{lemma-pullback-composition}
Let $f : X \to Y$ and $g : Y \to Z$ be continuous maps
of topological spaces. The functors $(g \circ f)^{-1}$
and $f^{-1} \circ g^{-1}$ are canonically isomorphic.
Similarly $(g \circ f)_p \cong f_p \circ g_p$ on
presheaves.
\end{lemma}

\begin{proof}
To see this
use that adjoint functors are unique up to unique isomorphism,
and Lemma \ref{lemma-pushforward-composition}.
\end{proof}

\begin{definition}
\label{definition-f-map}
Let $f : X \to Y$ be a continuous map.
Let $\mathcal{F}$ be a sheaf of sets on $X$ and
let $\mathcal{G}$ be a sheaf of sets on $Y$.
An {\it $f$-map $\xi : \mathcal{G} \to \mathcal{F}$}
is a collection of maps
$\xi_V : \mathcal{G}(V) \to \mathcal{F}(f^{-1}(V))$
indexed by open subsets $V \subset Y$ such that
$$
\xymatrix{
\mathcal{G}(V) \ar[r]_{\xi_V} \ar[d]_{\text{restriction of }\mathcal{G}} &
\mathcal{F}(f^{-1}V) \ar[d]^{\text{restriction of }\mathcal{F}} \\
\mathcal{G}(V') \ar[r]^{\xi_{V'}} &
\mathcal{F}(f^{-1}V')
}
$$
commutes for all $V' \subset V \subset Y$ open.
\end{definition}

\noindent
In the literature we sometimes find this defined alternatively
as in part (4) of Lemma \ref{lemma-f-map} but as the lemma shows there
is really no difference.

\begin{lemma}
\label{lemma-f-map}
Let $f : X \to Y$ be a continuous map.
There are bijections between the following four sets
\begin{enumerate}
\item the set of maps $\mathcal{G} \to f_*\mathcal{F}$,
\item the set of maps $f^{-1}\mathcal{G} \to \mathcal{F}$,
\item the set of $f$-maps $\xi : \mathcal{G} \to \mathcal{F}$, and
\item the set of all collections of maps
$\xi_{U, V} : \mathcal{G}(V) \to \mathcal{F}(U)$ for all
$U \subset X$ and $V \subset Y$ open such that $f(U) \subset V$
compatible with all restriction maps,
\end{enumerate}
functorially in $\mathcal{F} \in \Sh(X)$ and $\mathcal{G} \in \Sh(Y)$.
\end{lemma}

\begin{proof}
A map of sheaves $a : \mathcal{G} \to f_*\mathcal{F}$
is by definition a rule which to each open $V$ of $Y$ assigns
a map $a_V : \mathcal{G}(V) \to f_*\mathcal{F}(V)$
and we have $f_*\mathcal{F}(V) = \mathcal{F}(f^{-1}(V))$.
Thus at least the "data" corresponds exactly to what you need
for an $f$-map $\xi$ from $\mathcal{G}$ to $\mathcal{F}$.
To show that the sets (1) and (3) are in bijection we observe that
$a$ is a map of sheaves if and only if corresponding family of maps $a_V$
satisfy the condition in Definition \ref{definition-f-map}.

\medskip\noindent
Recall that $f^{-1}\mathcal{G}$ is the sheafification of $f_p\mathcal{G}$.
By the universal property of sheafification a map of sheaves
$b : f^{-1}\mathcal{G} \to \mathcal{F}$ is the same thing as a map of
presheaves $b_p : f_p\mathcal{G} \to \mathcal{F}$ where $f_p$ is the
functor defined earlier in the section. To give such a map $b_p$
you need to specify for each open $U$ of $X$ a map
$$
b_{p, U} :
\colim_{f(U) \subset V} \mathcal{G}(V)
\longrightarrow
\mathcal{F}(U)
$$
compatible with restriction mappings. We may and do view $b_{p, U}$
as a collection of maps $b_{p, U, V} : \mathcal{G}(V) \to \mathcal{F}(U)$
for all $V$ open in $Y$ with $f(U) \subset V$. These maps have to be
compatible with all possible restriction mappings you can think of.
In other words, we see that $b_p$ corresponds to a collection of maps
as in (4). Of course, conversely such a collection defines a map
$b_p$ and in turn a map $b : f^{-1}\mathcal{G} \to \mathcal{F}$.

\medskip\noindent
To finish the proof of the lemma you have to show that by
"forgetting structure" the rule that to a collection $\xi_{U, V}$
as in (4) associates the $f$-map $\xi$ with $\xi_V = \xi_{f^{-1}(V), V}$
is bijective. To do this, if $\xi$ is a usual $f$-map then we just
define $\tilde \xi_{U, V}$ to be the composition of
$\xi_V : \mathcal{G}(V) \to \mathcal{F}(f^{-1}(V))$
by the restruction map $\mathcal{F}(f^{-1}(V)) \to \mathcal{F}(U)$
which makes sense exactly because $f(U) \subset V$, i.e.,
$U \subset f^{-1}(V)$. This finishes the proof.
\end{proof}

\noindent
It is sometimes convenient to think about $f$-maps
instead of maps between sheaves either on $X$ or on $Y$.
We define composition of $f$-maps as follows.

\begin{definition}
\label{definition-composition-f-maps}
Suppose that $f : X \to Y$ and $g : Y \to Z$ are continuous
maps of topological spaces. Suppose that $\mathcal{F}$ is
a sheaf on $X$, $\mathcal{G}$ is a sheaf on $Y$, and
$\mathcal{H}$ is a sheaf on $Z$.
Let $\varphi : \mathcal{G} \to \mathcal{F}$ be an $f$-map.
Let $\psi : \mathcal{H} \to \mathcal{G}$ be an $g$-map.
The {\it composition of $\varphi$ and $\psi$} is the
$(g \circ f)$-map $\varphi \circ \psi$ defined
by the commutativity of the diagrams
$$
\xymatrix{
\mathcal{H}(W) \ar[rr]_{(\varphi \circ \psi)_W}
\ar[rd]_{\psi_W} & &
\mathcal{F}(f^{-1}g^{-1}W) \\
&
\mathcal{G}(g^{-1}W)
\ar[ru]_{\varphi_{g^{-1}W}}
}
$$
\end{definition}

\noindent
We leave it to the reader to verify that this works.
Another way to think about this is to think of
$\varphi \circ \psi$ as the composition
$$
\mathcal{H}
\xrightarrow{\psi}
g_*\mathcal{G}
\xrightarrow{g_*\varphi}
g_* f_* \mathcal{F} = (g \circ f)_* \mathcal{F}
$$
Now, doesn't it seem that thinking about $f$-maps is somehow
easier?

\medskip\noindent
Finally, given a continuous map $f : X \to Y$, and an
$f$-map $\varphi : \mathcal{G} \to \mathcal{F}$ there is
a natural map on stalks
$$
\varphi_x : \mathcal{G}_{f(x)} \longrightarrow \mathcal{F}_x
$$
for all $x \in X$. The image of a representative $(V, s)$
of an element in $\mathcal{G}_{f(x)}$ is mapped to the
element in $\mathcal{F}_x$ with representative $(f^{-1}V,
\varphi_V(s))$. We leave it to the reader to see that this
is well defined. Another way to state it is that it is the
unique map such that all diagrams
$$
\xymatrix{
\mathcal{F}(f^{-1}V) \ar[r] &
\mathcal{F}_x \\
\mathcal{G}(V) \ar[r] \ar[u]^{\varphi_V} &
\mathcal{G}_{f(x)} \ar[u]^{\varphi_x}
}
$$
(for $f(x) \in V \subset Y$ open) commute.

\begin{lemma}
\label{lemma-compose-f-maps-stalks}
Suppose that $f : X \to Y$ and $g : Y \to Z$ are continuous
maps of topological spaces. Suppose that $\mathcal{F}$ is
a sheaf on $X$, $\mathcal{G}$ is a sheaf on $Y$, and
$\mathcal{H}$ is a sheaf on $Z$.
Let $\varphi : \mathcal{G} \to \mathcal{F}$ be an $f$-map.
Let $\psi : \mathcal{H} \to \mathcal{G}$ be an $g$-map.
Let $x \in X$ be a point. The map on stalks
$(\varphi \circ \psi)_x : \mathcal{H}_{g(f(x))}
\to \mathcal{F}_x$ is the composition
$$
\mathcal{H}_{g(f(x))}
\xrightarrow{\psi_{f(x)}}
\mathcal{G}_{f(x)}
\xrightarrow{\varphi_x}
\mathcal{F}_x
$$
\end{lemma}

\begin{proof}
Immediate from Definition \ref{definition-composition-f-maps}
and the definition of the map on stalks above.
\end{proof}



